\documentclass[10pt, letterpaper]{article}

\usepackage[margin=0.85in]{geometry}
\usepackage{setspace}
\usepackage{titlesec}
\usepackage{fancyhdr}
\usepackage{microtype}
\usepackage{enumitem}
\usepackage{booktabs}
\usepackage{tabularx}
\usepackage{array}
\usepackage[hidelinks]{hyperref}

\setstretch{1.03}
\setlength{\parindent}{0.3in}
\setlength{\parskip}{1pt}
\setlist{leftmargin=*, itemsep=1pt, topsep=2pt}

\titleformat{\section}{\normalfont\normalsize\bfseries}{\thesection.}{0.5em}{}
\titlespacing*{\section}{0pt}{7pt}{3pt}

\pagestyle{fancy}
\fancyhf{}
\lhead{\small Md Kamruzzaman Kamrul}
\rhead{\small EDCI 59100-016 | Summer 2026}
\cfoot{\thepage}
\renewcommand{\headrulewidth}{0.4pt}
\setlength{\headheight}{14pt}

\newcolumntype{Y}{>{\raggedright\arraybackslash}X}

\begin{document}

\begin{center}
  {\Large\bfseries Final Literature Review Plan}\\[0.35em]
  {\normalsize Md Kamruzzaman Kamrul}\\
  {\normalsize EDCI 59100-016 | Summer 2026}
\end{center}

\section{Topic}

My review focuses on how vision-language models (VLMs) and related multimodal AI systems are being used for safety monitoring and activity tracking on active construction sites. In this plan, \textit{spatiotemporal reasoning} means understanding what is happening, where it is happening, and how it changes over time. For construction, that includes whether a worker is wearing PPE, whether equipment is idle or active, whether a task sequence is progressing, and whether a visual scene matches a safety rule or work report.

\section{Research Questions and Purpose}

The purpose of the review is to map how this literature is developing across model architectures, data sources, construction tasks, and evaluation practices. I am not trying to prove that one model is universally best. I am identifying what has been studied, what kinds of evidence exist, and where the literature is still thin or uneven.

\begin{enumerate}[label=\textbf{RQ\arabic*:}]
  \item What VLM or multimodal architectures are being used for construction-site safety monitoring and activity tracking?
  \item How are visual, textual, temporal, and geometric data combined to support construction safety or progress decisions?
  \item What practical and methodological barriers appear when these models are used for real-time or field-based monitoring?
  \item How are these systems evaluated, and what metrics or benchmarks are used across studies?
\end{enumerate}

These questions fit a scoping review because they ask what evidence exists, how it is organized, and where gaps remain. Each question will also guide a data extraction category: model type, data modality, construction task, deployment setting, evaluation metric, and stated limitation.

\section{Selected Review Methodology}

I will use a scoping review, reported with PRISMA-ScR guidance. A scoping review is appropriate because this topic is still forming across construction management, computer vision, natural language processing, and robotics. The studies do not use one shared design or one shared outcome measure. Some test zero-shot PPE monitoring. Others examine image captioning, daily report generation, equipment task monitoring, robotic inspection, or semantic safety reasoning.

For that reason, the review will emphasize mapping and categorization. I will chart what kinds of models, tasks, datasets, and evaluation approaches appear in the literature. PRISMA-ScR will help keep the process transparent, but I will avoid treating the project like a narrow systematic review or meta-analysis. The main goal is not effect-size comparison. The main goal is to show the shape of the literature in a usable way.

\section{Search Strategy}

My first broad search was too general. Searching terms such as ``AI in construction'' produced more than 8,000 results and mixed together unrelated topics such as project management automation, BIM, scheduling, structural inspection, and general computer vision. That search helped me see that the review needed clearer boundaries and stronger Boolean logic.

The refined search strategy now uses three concept groups:

\begin{itemize}
  \item \textbf{Architecture terms:} ``vision-language model,'' VLM, CLIP, image captioning, VQA, large multimodal model, foundation model, generative AI, video-language.
  \item \textbf{Task terms:} spatiotemporal, temporal, video, dynamic, sequence, tracking, action recognition, activity recognition, progress monitoring.
  \item \textbf{Construction terms:} construction site, construction management, AEC, civil engineering, safety monitoring, PPE, productivity, digital twin.
\end{itemize}

The final search combined these groups with \textbf{AND} and used \textbf{OR} within each group. This worked better because it forced each result to connect model type, task, and construction context. Web of Science returned 769 records and Scopus returned 674 records, for 1,443 raw records. After filtering and 33 duplicate removals, 90 unique articles remained. Title and abstract screening removed 32 articles outside the active-construction scope, leaving 58. Full-text review produced 25 core or transitional studies. Backward snowballing added 31 baseline and transition studies, creating a working corpus of 56 sources.

I will continue refining the search by checking whether important included papers use terms that my current search misses. For example, some construction papers use ``image captioning'' or ``semantic scene understanding'' instead of VLM. Others use ``activity monitoring'' rather than spatiotemporal reasoning. I will record these changes rather than silently changing the search.

\section{Databases and Sources}

The main databases are Web of Science Core Collection and Scopus. These are appropriate because the topic crosses civil engineering, construction management, computer science, automation, and applied AI. Web of Science gives strong coverage of high-impact indexed journals. Scopus adds broader engineering and interdisciplinary coverage.

I will use backward snowballing from included papers to find earlier baseline studies. This is important because some necessary background papers predate current VLM terminology. For example, RFID tracking, wearable sensors, CNN-based PPE detection, and image-based construction analysis are not the center of the review, but they help explain why current multimodal methods matter. Google Scholar may be used for citation tracing, but not as a primary database because its results are less reproducible.

\section{Inclusion and Exclusion Criteria}

The criteria were revised after screening. I now separate core VLM studies from baseline studies instead of treating every relevant AI paper the same way.

\begin{center}
\small
\renewcommand{\arraystretch}{1.2}
\begin{tabularx}{\textwidth}{p{2.4cm}YY}
\toprule
\textbf{Category} & \textbf{Include} & \textbf{Exclude} \\
\midrule
Timeframe & 2019--2026 for core VLM and multimodal AI studies; 2010--2026 for historical baselines. & Older studies unless needed as a clearly justified background reference. \\
Source type & Peer-reviewed journal articles; selected high-value review articles for context. & Editorials, theses, book chapters, industry white papers, and non-English sources. \\
Technology & VLMs, CLIP-based methods, VQA, image captioning, large multimodal models, hybrid vision-language systems, and selected transition methods. & Pure text-only NLP or pure object detection unless used as a baseline or bridge to multimodal reasoning. \\
Context & Active construction sites, construction safety, PPE monitoring, equipment/task tracking, productivity, progress reporting, or site inspection. & Post-construction facility management, general robotics outside construction, or structural health monitoring with no active-site relevance. \\
\bottomrule
\end{tabularx}
\end{center}

Screening decisions will use a simple Keep, Maybe, Exclude structure. ``Keep'' means the study directly supports the VLM or multimodal construction focus. ``Maybe'' means the study is useful as context or a baseline but is not central. ``Exclude'' means it fails the construction context, source type, or technology criteria.

\section{Emerging Patterns}

Two early patterns are visible in the source matrix. First, the field is moving from detection toward semantic interpretation. Older studies focus on cameras, RFID, wearable sensors, CNN detection, or activity recognition. Newer studies add language as a reasoning layer: captioning site images, matching visual scenes to safety rules, asking visual questions in augmented reality, or using MLLMs to interpret equipment activity.

Second, temporal reasoning is still a weak point. Several papers can identify objects, PPE, or activities in still images or short clips, but fewer systems handle change over time with enough reliability for field deployment. The Jeoung et al. equipment-monitoring study is useful here because it separates spatial, task, and temporal event questions. That distinction will help me chart whether studies are truly spatiotemporal or only visual-semantic.

\section{Reflection on Revision}

The largest change in my plan is that I simplified the methodology description. Earlier versions leaned too much toward systematic-review language. The plan is now more clearly framed as a scoping review: chart the literature, categorize heterogeneous studies, and identify gaps. I also added a plain definition of spatiotemporal reasoning because the term can sound more technical than it needs to be.

The search process also changed my criteria. At first, I expected the review to focus only on VLMs. The actual search showed that the construction literature is still transitioning. Some important studies use image captioning, CLIP, VQA, or hybrid vision-language reasoning without always using the term ``VLM.'' I also found that older baseline studies are necessary for explaining the shift from sensing and object detection toward semantic, multimodal interpretation.

Going forward, I will keep the extraction process manageable by charting every included study with the same categories: article, model or method, data modality, construction task, temporal component, evaluation metric, deployment context, and limitation. This should make the final synthesis easier to follow and prevent the review from becoming only a long list of papers.

\end{document}
